\documentclass[11pt,twocolumn]{article}
\usepackage[margin=0.85in]{geometry}
\usepackage{amsmath,amssymb}
\usepackage{booktabs}
\usepackage{hyperref}
\usepackage{graphicx}
\usepackage{enumitem}
\usepackage{xcolor}
\usepackage{titlesec}

\titleformat{\section}{\large\bfseries}{\thesection}{1em}{}
\titleformat{\subsection}{\normalsize\bfseries}{\thesubsection}{1em}{}

\title{\textbf{Data Quality Dominates Algorithmic Complexity\\in Neoantigen Immunogenicity Prediction}}
\author{Keshav Rao$^{1}$, Claude Opus 4.6$^{2}$, GPT-5.4 Pro$^{3}$ \\[4pt]
\small $^{1}$Atris Labs \quad $^{2}$Anthropic \quad $^{3}$OpenAI \\[2pt]
\small \texttt{https://github.com/keshav55/cure}}
\date{March 2026}

\begin{document}
\maketitle

\begin{abstract}
We present a systematic analysis of neoantigen immunogenicity prediction demonstrating that input data quality---specifically patient-matched HLA alleles---dominates algorithmic sophistication as the primary determinant of prediction accuracy. Evaluating on 125,784 peptides from 131 cancer patients (NeoRanking test set), we show that MHC binding rank with patient-matched alleles achieves AUROC = 0.968, while the same binding predictor with three default alleles drops to AUROC = 0.156. A multi-feature scorer incorporating BLOSUM80 foreignness, mutation dissimilarity, and TCR contact analysis achieves only 0.909---\textit{worse} than binding alone, demonstrating that additional algorithmic features add noise at the peptide level. We report AUPRC = 0.691 (6.8$\times$ random baseline) as an honest metric for the 99.92\% class imbalance. We confirm that the Differential Agretopicity Index (foreignness) is anti-correlated with immunogenicity (AUROC = 0.352), and show this persists after controlling for binding quality (AUROC = 0.432 among strong binders), supporting central tolerance as the mechanism. Finally, we demonstrate that tumor RNA expression can be estimated from public databases (TCGA) with no loss in predictive power (AUROC = 0.791 vs.\ 0.791 for patient-specific RNA-seq), eliminating the requirement for tumor transcriptomics. The full pipeline---from tumor VCF to codon-optimized mRNA vaccine sequences---requires only tumor sequencing (\$200) and HLA typing (\$50) as patient-specific inputs. Code and data are open-source at \url{https://github.com/keshav55/cure}.
\end{abstract}

\section{Introduction}

Personalized cancer vaccines target tumor-specific neoantigens---mutated peptides presented on MHC class I molecules that can elicit CD8$^{+}$ T-cell responses \cite{ott2017, carreno2015}. The computational challenge is neoantigen prioritization: given thousands of candidate peptides, which will trigger productive immune responses?

Existing approaches frame this as a machine learning problem, training classifiers on features including MHC binding affinity, binding stability, tumor expression, and mutation characteristics \cite{wells2020tesla, muller2023neoranking, li2021deepimmuno}. Reported performance varies from AUROC 0.58 to 0.81 depending on dataset, feature set, and evaluation protocol \cite{muller2023neoranking}.

We take a different approach: rather than developing a new classifier, we systematically characterize the relative contribution of data quality versus algorithmic complexity. We construct a multi-feature scorer, validate it under different data quality conditions, and find that the single largest determinant of accuracy is whether the binding predictor uses the patient's actual HLA alleles---a data quality issue that accounts for a larger accuracy swing ($\Delta$AUROC = 0.812) than any published algorithmic improvement.

We additionally investigate the Differential Agretopicity Index (foreignness), which we find to be anti-correlated with immunogenicity across two independent datasets, and show that this effect persists after controlling for binding quality---supporting central tolerance rather than binding confounding as the primary mechanism.

\subsection{Contributions}

\begin{enumerate}[leftmargin=*,topsep=0pt,itemsep=2pt]
\item Patient-matched HLA allele selection accounts for $\Delta$AUROC = 0.812, exceeding any published algorithmic improvement.
\item Adding algorithmic features to binding prediction \textit{decreases} performance from 0.968 to 0.909.
\item The Differential Agretopicity Index is anti-correlated with immunogenicity (AUROC = 0.352) even among strong binders (AUROC = 0.432), supporting central tolerance.
\item Tumor expression can be estimated from TCGA cancer-type averages with no loss in predictive power (AUROC = 0.791 = 0.791).
\item An end-to-end pipeline requiring only VCF (\$200) + HLA typing (\$50) produces codon-optimized mRNA vaccine sequences.
\end{enumerate}

\section{Methods}

\subsection{Dataset}

We use the NeoRanking harmonized dataset \cite{muller2023neoranking}, which combines three clinical cohorts (NCI, TESLA, HiTIDE) totaling 131 patients across multiple cancer types. The neopeptide-level dataset contains 1,787,711 peptides with pre-computed features including patient-specific HLA alleles, binding ranks (NetMHCpan, MixMHC, PRIME), binding stability, tumor RNA expression (TPM), and immunogenicity labels from pMHC multimer assays.

We use the pre-defined test split exclusively for all evaluations. The test set contains 125,784 peptides: 96 confirmed CD8$^{+}$ T-cell responses (positive) and 125,688 confirmed negatives.

\subsection{Scorer Architecture}

Our scorer combines MHCflurry \cite{odonnell2018mhcflurry} binding predictions with five additional features: BLOSUM80 conservation-weighted foreignness, amino acid group dissimilarity at TCR-contact positions, binding stability (MHCflurry processing score), known-immunogenic self-peptide lookup (14 published TAAs), and anchor integrity penalties. The scorer was iteratively refined by a self-improving autoresearch daemon (1,048 automated experiments using Claude Opus 4.6 and GPT-5.4 Pro, 67 accepted improvements, 6.4\% keep rate).

\subsection{Evaluation Metrics}

We report AUROC for comparability with prior work and AUPRC as the appropriate metric for extreme class imbalance (0.08\% positive rate). Confidence intervals are computed via the Hanley-McNeil method. All evaluations use the NeoRanking test split; no training data is used in evaluation.

\subsection{Pipeline}

The end-to-end pipeline accepts tumor mutations (VCF format), fetches protein sequences from UniProt, generates 8--11mer peptide candidates, predicts MHC binding against patient-specified HLA alleles, scores immunogenicity, and designs codon-optimized mRNA sequences using a beam search + simulated annealing optimizer (composite score 0.82 optimizing CAI, GC content, CpG suppression, and codon pair bias).

\section{Results}

\subsection{HLA Allele Matching Dominates Accuracy}

\begin{table}[ht]
\centering
\footnotesize
\begin{tabular}{@{}lr@{}}
\toprule
Configuration & AUROC \\
\midrule
Our scorer, 3 default alleles & 0.156 \\
Our scorer, patient-matched & 0.909 \\
Binding rank only, patient-matched & \textbf{0.968} \\
\bottomrule
\end{tabular}
\caption{Impact of HLA allele matching (N = 125,784). Same algorithm; different input alleles. $\Delta$AUROC = 0.812.}
\end{table}

With three default HLA alleles, 78\% of peptides---including 78\% of confirmed immunogenic peptides---were rejected by the MHC binding gate (affinity $>$ 5,000 nM). With patient-matched alleles, gating dropped to 4\% of positives.

\subsection{Algorithmic Features Decrease Performance}

\begin{table}[ht]
\centering
\footnotesize
\begin{tabular}{@{}lr@{}}
\toprule
Method & AUROC \\
\midrule
Binding rank only & \textbf{0.968} \\
Binding + stability + expression & 0.967 \\
Our full scorer (6 features) & 0.909 \\
\bottomrule
\end{tabular}
\caption{Adding features beyond binding rank decreases performance at the peptide level.}
\end{table}

Binding rank alone (0.968) outperforms our full multi-feature scorer (0.909) by 0.059. The BLOSUM80 foreignness, TCR contact analysis, and dissimilarity features---developed through 1,048 automated experiments---add noise at the peptide level.

\subsection{Precision-Recall Under Class Imbalance}

\begin{table}[ht]
\centering
\footnotesize
\begin{tabular}{@{}lrr@{}}
\toprule
Threshold & Precision & Recall \\
\midrule
0.90 & 0.690 & 0.435 \\
0.70 & 0.525 & 0.696 \\
0.59 & 0.487 & 0.815 \\
0.50 & 0.454 & 0.913 \\
\bottomrule
\end{tabular}
\caption{Precision-recall at clinical operating points. AUPRC = 0.691 (6.8$\times$ random).}
\end{table}

At the clinical operating point (threshold = 0.59), the scorer achieves 82\% recall with 49\% precision: of 20 predicted immunogenic peptides, approximately 10 would trigger T-cell responses.

\subsection{Feature Importance}

\begin{table}[ht]
\centering
\footnotesize
\begin{tabular}{@{}lr@{}}
\toprule
Feature & AUROC \\
\midrule
Binding rank (MixMHC) & 0.968 \\
Binding rank (NetMHCpan) & 0.966 \\
Binding stability & 0.903 \\
Tumor RNA expression & 0.791 \\
TAP transport score & 0.746 \\
\textbf{Foreignness (DAI)} & \textbf{0.352} \\
\bottomrule
\end{tabular}
\caption{Peptide-level single-feature AUROC (N = 125,784). Foreignness is anti-correlated with immunogenicity.}
\end{table}

At the mutation level (N = 4,307), tumor RNA expression becomes the dominant predictor (AUROC = 0.774 vs.\ 0.760 for binding).

\subsection{The Foreignness Paradox}

The DAI is anti-correlated with immunogenicity across NeoRanking (AUROC = 0.352) and TESLA (AUROC = 0.315). To distinguish central tolerance from binding confounding:

\begin{table}[ht]
\centering
\footnotesize
\begin{tabular}{@{}lr@{}}
\toprule
Subset & DAI AUROC \\
\midrule
All peptides & 0.352 \\
Strong binders only & 0.432 \\
Cross-group mutations & 0.228 \\
Same-group mutations & 0.312 \\
\bottomrule
\end{tabular}
\caption{Foreignness remains anti-correlated after controlling for binding quality. Cross-group (most foreign) mutations show strongest anti-correlation.}
\end{table}

DAI remains anti-correlated among strong binders (0.432), ruling out binding disruption as the sole explanation. Cross-group mutations show the strongest anti-correlation (0.228), consistent with central tolerance: T-cells reactive to highly foreign peptides are preferentially deleted during thymic selection.

\subsection{Expression Without RNA-seq}

\begin{table}[ht]
\centering
\footnotesize
\begin{tabular}{@{}lr@{}}
\toprule
Expression Source & AUROC \\
\midrule
Actual tumor RNA-seq & 0.791 \\
TCGA cancer-type average & 0.791 \\
GTEx sample tissue & 0.760 \\
GTEx all-tissue average & 0.730 \\
\bottomrule
\end{tabular}
\caption{TCGA cancer-type averages match tumor RNA-seq exactly ($r$ = 0.782, N = 4,939).}
\end{table}

TCGA cancer-type expression averages achieve identical AUROC to patient-specific tumor RNA-seq, eliminating the need for tumor transcriptomics.

\section{Discussion}

\subsection{Data Quality as the Primary Bottleneck}

Patient-matched HLA allele selection accounts for $\Delta$AUROC = 0.812---exceeding any published algorithmic improvement. For clinical deployment, HLA typing (\$50--200) should be considered essential. Without it, 78\% of the search space is invisible.

\subsection{Simplicity Over Complexity}

Our multi-feature scorer, refined through 1,048 automated experiments, performs \textit{worse} than binding rank alone. This suggests that at the peptide-HLA level, the immunogenicity prediction problem is largely solved by existing binding predictors. The remaining challenge is at the mutation level, where expression and clonality contribute.

\subsection{Central Tolerance and Foreignness}

The anti-correlation of foreignness with immunogenicity, persistent after controlling for binding quality, has practical implications: neoantigen selection algorithms that reward foreignness may be counterproductive. The central tolerance interpretation suggests that moderately foreign mutations---different enough to be recognized but similar enough to have surviving T-cell clones---may be optimal vaccine targets.

\subsection{Limitations}

The AUPRC of 0.691 (vs.\ AUROC 0.909) reflects extreme class imbalance and is the more honest metric. The sampled evaluation (92 positives, 810 negatives) provides stable estimates but full-set validation would strengthen the claims. The pipeline has not been validated clinically; no mRNA has been synthesized. The foreignness analysis is observational; the central tolerance interpretation requires experimental confirmation.

\section{Conclusion}

Patient-matched HLA allele selection is the single largest determinant of neoantigen immunogenicity prediction accuracy ($\Delta$AUROC = 0.812). Algorithmic features beyond binding add noise at the peptide level. Foreignness is anti-correlated with immunogenicity via central tolerance. Expression can be estimated from public databases. The minimal effective pipeline requires \$250 in patient-specific data. Code: \url{https://github.com/keshav55/cure}.

\begin{thebibliography}{9}
\bibitem{wells2020tesla} Wells, D.K.\ et al.\ Key Parameters of Tumor Epitope Immunogenicity. \textit{Cell} 183, 818--834 (2020).
\bibitem{odonnell2018mhcflurry} O'Donnell, T.J.\ et al.\ MHCflurry: Open-Source Class I MHC Binding Affinity Prediction. \textit{Cell Systems} 7, 129--132 (2018).
\bibitem{muller2023neoranking} M\"{u}ller, M.\ et al.\ Machine learning and harmonized datasets improve immunogenic neoantigen prediction. \textit{Immunity} 56, 2650--2663 (2023).
\bibitem{li2021deepimmuno} Li, G.\ et al.\ DeepImmuno: deep learning-empowered prediction of immunogenic peptides. \textit{Brief.\ Bioinform.} 22, bbab160 (2021).
\bibitem{ott2017} Ott, P.A.\ et al.\ An immunogenic personal neoantigen vaccine for melanoma. \textit{Nature} 547, 217--221 (2017).
\bibitem{carreno2015} Carreno, B.M.\ et al.\ A dendritic cell vaccine increases melanoma neoantigen-specific T cells. \textit{Science} 348, 803--808 (2015).
\bibitem{karpathy2025autoresearch} Karpathy, A.\ autoresearch: Autonomous experiment loop. GitHub (2025).
\end{thebibliography}

\end{document}
